\section{Related work}

We aim to re-render a video from new viewpoints with temporal control, a task closely related to Novel View Synthesis (NVS) from monocular video inputs.




\myparagraph{Video-based NVS.}
Prior video-based NVS methods can be broadly characterized along two axes: (i) whether they target static or dynamic scenes, and (ii) whether they incorporate explicit 3D geometry in the generation pipeline.


For static scenes, geometry-based methods reconstruct scene geometry from the input frames and use diffusion models to complete or hallucinate regions that are unseen under new viewpoints \cite{yu2024viewcrafter,Ren2025,Jeong2025,wu2025video}.
Although these approaches achieve high rendering quality, they rely on heavy 3D preprocessing.
Geometry-free approaches \cite{Bai2025,Zhou2025,vanhoorick2024gcd} bypass explicit geometry and directly condition the diffusion process on observed views and camera poses to synthesize new viewpoints.


For dynamic scenes, inpainting-based methods such as TrajectoryCrafter \cite{YU2025}, ReCapture \cite{Zhang2024}, and Reangle \cite{Jeong2025} also adopt warp-and-inpaint pipelines, while GEN3C \cite{Ren2025} extends this with an evolving 3D cache and EPiC \cite{zun2025epic} improves efficiency via a lightweight ControlNet framework.
Geometry-free dynamic models \cite{Bai2025,bai2024syncammaster,vanhoorick2024gcd, wang20254realv2, Wang_2025_CVPR} instead learn camera-conditioned generation from multi-view or 4D datasets (\eg, Kubric-4D \cite{greff2021kubric}), enabling smoother and more stable NVS with minimal 3D inductive bias.
Proprietary systems like Genie 3 \cite{parker-holder2025genie3} further demonstrate real-time, continuous camera control in dynamic scenes, underscoring the potential of video diffusion models for interactive viewpoint manipulation.


\myparagraph{Disentangling Space and Time.}
Despite great progress in camera controllability (space), 
the methods discussed above do not address temporal control (time).
Meanwhile, disentangling spatial and temporal factors has become a central focus in 4D scene generation, recently advanced through diffusion-based models.
\fourdim \cite{watson2025controllingspacetimediffusion} introduces a Masked FiLM mechanism that defaults to identity transformations when conditioning signals (e.g., camera pose or time) are absent, enabling unified representations across both static and dynamic data through multi-modal supervision. Similarly, \catfourd \cite{Wu2024} leverages multi-view images to conduct 4D dynamic reconstruction to achieve space–time disentanglement but remains constrained by its reliance on explicit 4D reconstruction pipelines, which limits scalability and controllability. 
In contrast, our approach builds upon text-to-video diffusion models and introduces a new temporal embeddings module and refined camera conditioning to achieve fully controllable 4D generative reconstruction. 

